\noindent\begin{tabularx}{\linewidth}{@{}L{4.0cm}Y@{}}
\toprule
Statement & Opened source and scope\\
\midrule
Two machines & Coffman--Graham \cite[Theorem~1, p.~203]{CG}; Fujii--Kasami--Ninomiya \cite[abstract]{FKN} with the correction \cite[p.~141]{FKNe}.\\[2pt]
Machine count in the input & Ullman's unit-time NP-completeness result \cite[Theorem~1, p.~389]{U}; the fixed-count question is discussed on p.~392.\\[2pt]
Three machines & OpenAI \cite[Theorem~1.1]{P}. The earlier question is OPEN8 \cite[Appendix A13, p.~287]{GJ}. Giraudin's subsequent statement has input exponent $276$ \cite[Theorem~1]{G}.\\[2pt]
Four machines & drbh's commit \texttt{bbd3013}, \cite[\texttt{Model.lean}, \texttt{ReleaseTheorem}; \texttt{Main.lean}, \texttt{main\_theorem}]{D}, contains the capacity-four solver and theorem code.\\[2pt]
Every fixed $m\geq3$ & Jansen \cite[Theorem~1.1]{J} states exact decision and optimization; \cite[Lemma~7.15]{J} gives node/cut leaf budgets $4m+5$ and $m+3$.\\[2pt]
Exact subexponential bound & Nederlof--Swennenhuis--W\k{e}grzycki \cite[Theorem~1.1]{N}: $(1+n/m)^{O(\sqrt{nm})}$.\\[2pt]
Approximation scheme & Das--Wiese \cite[Theorem~1]{A}: $(1+\varepsilon)$ approximation with the parameter-dependent quasipolynomial bound stated there.\\[2pt]
This post & Theorems A and B are proved here with $12m+25$ leaves and the derived exponent in (1.1). No novelty or priority is claimed.\\
\bottomrule
\end{tabularx}

\smallskip
\noindent\textbf{Relation of the guarantees.}
Lemma 2.2 gives integer makespans, and a topological schedule gives $\mathrm{OPT}\leq n$.
Using $\varepsilon=1/(n+1)$ in the guarantee of \cite[Theorem~1]{A} yields a feasible integer makespan strictly below $\mathrm{OPT}+1$, hence exactly $\mathrm{OPT}$.
The source's time bound depends on this varying $\varepsilon$; (1.1) is derived directly from the finite search in this post.
For fixed $m$, the displayed bound of \cite[Theorem~1.1]{N} has the form $\exp(O_m(\sqrt n\log(n+1)))$, by taking its logarithm. This is the stated upper bound, with no claim of a corresponding lower bound on that algorithm.

\smallskip
\noindent\textbf{Credit and search.}
The construction follows OpenAI \cite[Sections 2--9]{P} and the course \cite[lessons 1--8]{L}.
The exact scope of drbh's four-machine code is credited in the source record; its Lean build was not replayed in this audit.
The search of 15 November 2026 covered arXiv, STOC, FOCS, SODA, ICALP, IPCO, ESA, APPROX, authors' pages, GitHub, the general web and Hexagon, including citations and title searches for [P]. It found [D], [J] and [G]; no fixed-$m\geq4$ hardness result was located. Every query and outcome is in \texttt{Q03\_DEPOSIT\_literature.md}.
